\section{An observed ``no'' is not a diagnosis}

The word \emph{refusal} describes what a user sees, not why it happened. A product may explicitly decline, repeatedly ask for more direction, stop after a draft, or fail silently before the external effect. These episodes feel similar to the customer because the task remains undone. They are analytically different. A useful research program must distinguish a model that cannot choose the next action from a product that lacks a tool, a session that lost context, an account without delegated scope, and a task correctly barred by policy or platform rule. Only some of these are aftermarket opportunities.

We propose recording a \emph{stop event} with four fields: the requested terminal state, the product's last observable action, the first unmet precondition, and the human intervention required to proceed. The unmet precondition is diagnosed with evidence, not with a guess about hidden model states. For example, a model may correctly describe an API update but have no connected API. That is evidence for a product integration gap. If an API is present but the model chooses the wrong record, the evidence points toward planning or perception. If the owner has not authorized a write, the gap is authority. If the request is properly prohibited, the stop is a correct result. The labels should be allowed to coexist when an episode has multiple causes.

Figure~\ref{fig:refusal} is a short decision tree for triage. It puts the proper boundary first because a model's apparent ability to generate instructions is irrelevant if the action is not permitted. For a permitted task, the harness can test a capability hypothesis without making safety refusal the target. The key experimental move is a paired trial with the same underlying model and an unchanged acceptance rule. A successful harnessed trial supports the narrower statement that the stock product's decline was not a hard inability under those conditions. A failed trial does not establish that no harness could work.

\begin{figure}[tbp]
\centering\begin{forest}
for tree={
  grow'=0,
  anchor=west,
  child anchor=west,
  parent anchor=east,
  edge={draw=ink!35,line width=.65pt},
  l sep=15mm,s sep=5mm,
  inner xsep=1pt,inner ysep=2pt,
  align=left,
  font=\small
}
[\textbf{Observed stop},text=ink
  [\textcolor{rust}{\textbf{Correct boundary}}\\Policy or platform disallows the effect, text=ink]
  [\textcolor{teal}{\textbf{Permitted task}},text=ink
    [Capability gap\\Model cannot choose a reliable next step]
    [Product gap\\Context{,} tool{,} session{,} or time is missing]
    [Authority gap\\Owner has not granted the specific action]
  ]
]
\end{forest}

\caption{A diagnostic tree for an observed stop. It is a coding scheme for evaluation, not a procedure for overriding legitimate restrictions. ``Permitted task'' still requires the owner and platform to authorize the particular route.}
\label{fig:refusal}
\end{figure}

Research on over-refusal establishes that benign requests can be rejected under some model and prompt conditions. XSTest intentionally contrasts safe prompts with unsafe ones; OR-Bench expands the study of difficult benign prompts \citep{xstest,orbench}. These are valuable calibration instruments, but their datasets and model versions do not measure the current rate at which 2026 consumer products reject commercially valuable workflows. A workflow study needs its own task population and must record product configuration, available tools, permission state, and date. Otherwise the label ``over-refusal'' can become a rhetorical substitute for measurement.

The default consumer product may also be right to decline a request because the requested action requires identity, consent, a platform-approved route, or a person’s decision. In that case the researcher should count a clean handoff or denial as the correct behavior. A harness can still create value by preparing context, preserving progress, and resuming after a human acts, but it should not conceal the person’s part. ``Without nudging'' means no repeated tactical coaching for ordinary task steps; it does not mean no required approval, no ownership check, or no human judgment where the task demands one.

We reserve \emph{recoverable stop} for a permitted task whose completion improves under a specified intervention. The intervention might be a better task specification, retrieval of relevant context, access to an allowed tool, durable state across sessions, or an independent verifier. These have different implementation costs and generalization prospects. The distinction matters commercially: prompt-only recovery may be copied quickly by a provider; a service-specific connector may be expensive to maintain; a verification contract may be a durable source of customer trust. By coding the cause and intervention separately, an experiment can discover where a specialist has real leverage instead of treating every ``I can't'' as a sales lead.
